% !TEX root = ../main.tex
\chapter{The Module System}
\label{ch:module_system}

A neural network is, structurally, a tree of stateful computation
units. Each unit carries parameters (the trainable tensors), holds
buffers (non-trainable state such as running statistics), composes
with other units (a model is built from layers, layers from
sub-layers), and produces outputs by applying a forward function to
inputs. \Lucid's \emph{module system} (\code{lucid.nn}) is the
machinery that organises this tree, names its parts, traverses it
for serialisation and parameter collection, dispatches the forward
function, and hooks into the autograd system for the backward.

The design follows \refframework's \code{torch.nn} closely: a
\code{Module} base class with \code{forward}, named parameters and
buffers, recursive child management, training-vs-eval mode, and
the standard hook surface. The interoperability matters because
the model zoo (\partref{part:model_zoo}) and a large fraction of
user code is structured around this idiom.

This chapter describes the \code{Module} class, the \code{Parameter}
type, the container modules (\code{ModuleList}, \code{ModuleDict},
\code{Sequential}, \code{ParameterList}, \code{ParameterDict}),
the state-dict serialisation contract, the training-mode flag, and
the hook system.

\section{Overview}
\label{sec:module_system:overview}

\code{lucid.nn.Module} is the base class for every neural-network
component in \Lucid. Subclassing requires only two things:

\begin{enumerate}
  \item Define \code{\_\_init\_\_(self, ...)} to construct sub-modules
    and parameters.
  \item Define \code{forward(self, ...)} that applies the layer's
    computation.
\end{enumerate}

Everything else --- parameter collection, recursive traversal,
state-dict, training mode, hooks --- the base class provides.

\begin{lstlisting}[style=lucid,language=LucidPython]
import lucid
import lucid.nn as nn

class TwoLayerMLP(nn.Module):
    def __init__(self, in_dim, hidden, out_dim):
        super().__init__()
        self.fc1 = nn.Linear(in_dim, hidden)
        self.fc2 = nn.Linear(hidden, out_dim)

    def forward(self, x):
        h = self.fc1(x).relu()
        return self.fc2(h)

model = TwoLayerMLP(784, 256, 10)
y = model(lucid.randn(32, 784))
\end{lstlisting}

The \code{model(x)} call dispatches to \code{model.forward(x)} via
the base class's \code{\_\_call\_\_}, which also fires registered
pre- and post-forward hooks (\secref{sec:module_system:hooks}).

\section{The Module Class}
\label{sec:module_system:module_class}

\code{Module} is implemented in \code{lucid/nn/module.py}. The
relevant state on each instance:

\begin{itemize}
  \item \code{\_parameters}: ordered dict of \code{Parameter} objects
    declared as direct attributes.
  \item \code{\_buffers}: ordered dict of non-parameter
    \code{Tensor} objects (e.g., running statistics).
  \item \code{\_modules}: ordered dict of child \code{Module}
    instances declared as direct attributes.
  \item \code{\_non\_persistent\_buffers\_set}: set of buffer names
    that should \emph{not} appear in \code{state\_dict()} (e.g.,
    transient scratch arrays).
  \item \code{training}: bool flag, default \code{True}, toggles
    behaviour of layers like dropout and batch norm.
  \item \code{\_forward\_hooks}, \code{\_backward\_hooks},
    \code{\_forward\_pre\_hooks},
    \code{\_state\_dict\_hooks},
    \code{\_load\_state\_dict\_pre\_hooks}: hook registries.
  \item \code{\_lucid\_signature}: cached \code{inspect.Signature}
    of the subclass's \code{forward} method, used by the tracer
    (\chref{ch:tracer}) to map positional arguments to graph
    inputs.
\end{itemize}

\subsection{The three dictionaries}
\label{ssec:module_system:three_dicts}

The separation into three distinct ordered dictionaries
(parameters, buffers, modules) is deliberate. The optimiser
sees only parameters; the device-move operation touches both
parameters and buffers; the recursive traversal walks modules.
A single bag of attributes would conflate these concerns. The
shape of the data also informs the API:

\begin{table}[t]
\centering
\small
\begin{tabular}{lcccc}
\toprule
Dict & Trainable & In state\_dict & In .parameters() & In .children() \\
\midrule
\code{\_parameters} & Yes & Yes & Yes & No \\
\code{\_buffers}    & No  & Yes (unless non-persistent) & No & No \\
\code{\_modules}    & ---  & Recurses & Recurses & Yes \\
\bottomrule
\end{tabular}
\caption[The three module dictionaries.]{Each dictionary serves
a distinct role. Recursive traversal of parameters and buffers
walks through \code{\_modules} children depth-first, deterministic
by declaration order. Non-persistent buffers (registered with
\code{persistent=False}) are visible at runtime but skipped by
serialisation.}
\label{tab:module_system:three_dicts}
\end{table}

The implementation: each public iterator (\code{parameters},
\code{buffers}, \code{children}, \code{modules}) is a thin
wrapper over \code{\_named\_members}, which recurses through
\code{\_modules} and yields from each sub-module's
corresponding dict. The walk is depth-first, in declaration
order, with cycle detection (a module can only be visited once
in case of accidental cyclic references via Python's reference
graph).

\subsection{Attribute interception}
\label{ssec:module_system:attribute_intercept}

Assigning a \code{Parameter}, a \code{Tensor} flagged as buffer, or
a \code{Module} as an attribute triggers \code{\_\_setattr\_\_} to
route the assignment to the appropriate dict:

\begin{lstlisting}[style=lucid,language=LucidPython]
class Layer(nn.Module):
    def __init__(self):
        super().__init__()
        self.weight = nn.Parameter(lucid.randn(10, 10))  # -> _parameters
        self.bias = nn.Parameter(lucid.zeros(10))         # -> _parameters
        self.activation = nn.ReLU()                      # -> _modules
        self.register_buffer('running_mean', lucid.zeros(10))  # -> _buffers
\end{lstlisting}

The interception is what makes the recursive traversal possible:
the base class can walk \code{\_modules} to find every sub-module,
walk \code{\_parameters} to find every parameter, and so on, all
without the user explicitly registering anything.

\subsection{Recursive parameter collection}
\label{ssec:module_system:recursive_params}

\code{module.parameters()} returns a generator over every
\code{Parameter} in the module tree (including those in sub-modules,
recursively). This is the standard input to optimisers:

\begin{lstlisting}[style=lucid,language=LucidPython]
optimizer = lucid.optim.Adam(model.parameters(), lr=1e-3)
\end{lstlisting}

The traversal uses depth-first order, deterministic by attribute
declaration order. \code{named\_parameters()} returns
\code{(name, parameter)} pairs with dotted names
(\code{fc1.weight}, \code{fc1.bias}, \code{fc2.weight}, etc.).

\subsection{Tree statistics}
\label{ssec:module_system:tree_stats}

\code{module.num\_parameters()} returns the total number of
trainable parameters across the tree. For a 100M-parameter model:

\begin{lstlisting}[style=lucid,language=LucidPython]
n = sum(p.numel() for p in model.parameters() if p.requires_grad)
\end{lstlisting}

The pattern is so common that the convenience function
\code{lucid.nn.count\_parameters(model)} is provided.

A richer statistics function, \code{lucid.nn.summary(model)},
walks the tree once and produces a structured report:

\begin{lstlisting}[style=lucid,language=LucidPython]
print(lucid.nn.summary(model))
# TwoLayerMLP                                    267,786 params
#   fc1: Linear(in=784, out=256)                 200,960
#     weight                  (256, 784)         200,704
#     bias                    (256,)                 256
#   fc2: Linear(in=256, out=10)                    2,570
#     weight                  (10, 256)            2,560
#     bias                    (10,)                   10
# Total trainable: 203,530   non-trainable: 0   buffers: 0
\end{lstlisting}

The summary respects the declaration order, prints buffer
counts separately from parameters, and flags any
\code{requires\_grad=False} entries (under the
``non-trainable'' tally). It is the standard model-inspection
tool used by the H12-compliant model-zoo doc-site introspector
(\chref{ch:zoo_architecture}).

\subsection{Walking the tree manually}
\label{ssec:module_system:walking_tree}

Three iterators expose the tree at different granularities:

\begin{lstlisting}[style=lucid,language=LucidPython]
# Direct children only (one level deep):
for name, child in model.named_children():
    print(name, type(child).__name__)
# fc1 Linear
# fc2 Linear

# Every descendant module (including self):
for name, mod in model.named_modules():
    print(name, type(mod).__name__)
# '' TwoLayerMLP
# 'fc1' Linear
# 'fc2' Linear

# Every parameter (recursive, named):
for name, p in model.named_parameters():
    print(name, p.shape)
# fc1.weight (256, 784)
# fc1.bias (256,)
# fc2.weight (10, 256)
# fc2.bias (10,)
\end{lstlisting}

The choice of iterator matters for performance: a model with
thousands of parameters re-iterated every step incurs a
non-trivial Python overhead. The convention is to materialise
\code{list(model.parameters())} once at optimiser setup; the
optimiser caches references and revisits them every step
without re-walking the tree.

\subsection{Memory cost of a module tree}
\label{ssec:module_system:memory_cost}

A module instance is a small Python object: a few hundred bytes
for the empty \code{Module}, plus the three dicts (a few KB).
The tensor storage is separate. For a 100M-parameter ResNet,
the tree itself is $\sim$\,80\,KB of Python objects; the tensor
storage is 400\,MB FP32 (the actual memory). The tree-walking
overhead is therefore negligible relative to the tensors.

\section{The Parameter Class}
\label{sec:module_system:parameter_class}

\code{lucid.nn.Parameter} is a subclass of \code{lucid.Tensor} with
two distinguishing properties:

\begin{itemize}
  \item \code{requires\_grad = True} by default.
  \item Assigned to a \code{Module} attribute, the base class
    automatically routes it into \code{\_parameters}.
\end{itemize}

Functionally, a \code{Parameter} is just a tensor; the type marker
exists only to signal ``this tensor is trainable and should be
tracked by the module system.''

\subsection{Initialisation}
\label{ssec:module_system:param_init}

A freshly constructed \code{Parameter} has random or zero values.
The standard pattern is to initialise via one of the schemes from
\chref{ch:initialization}:

\begin{lstlisting}[style=lucid,language=LucidPython]
class Linear(nn.Module):
    def __init__(self, in_features, out_features):
        super().__init__()
        self.weight = nn.Parameter(lucid.empty(out_features, in_features))
        self.bias = nn.Parameter(lucid.empty(out_features))
        self.reset_parameters()

    def reset_parameters(self):
        nn.init.kaiming_uniform_(self.weight, a=lucid.math.sqrt(5))
        bound = 1 / lucid.math.sqrt(self.weight.shape[1])
        nn.init.uniform_(self.bias, -bound, bound)
\end{lstlisting}

The \code{reset\_parameters} method is a convention: every standard
layer in \code{lucid.nn} defines it, and the user can call it to
re-initialise without constructing a new instance.

\subsection{Parameter vs buffer}
\label{ssec:module_system:param_vs_buffer}

A \emph{parameter} is trainable: it has \code{requires\_grad=True},
participates in autograd, and is updated by the optimiser. A
\emph{buffer} is non-trainable persistent state: it has
\code{requires\_grad=False}, does not participate in autograd, and
is saved/loaded alongside parameters in \code{state\_dict}.

The canonical buffer use case is BatchNorm's running statistics:
the running mean and variance are not gradient-updated but are
updated by the layer's forward pass and must persist across
training/eval and across checkpoint save/load.

\section{Container Modules}
\label{sec:module_system:containers}

\code{lucid.nn} provides several container modules for organising
sub-module collections.

\subsection{Sequential}
\label{ssec:module_system:sequential}

\code{Sequential} chains its sub-modules in order, passing each's
output to the next:

\begin{lstlisting}[style=lucid,language=LucidPython]
model = nn.Sequential(
    nn.Linear(784, 256),
    nn.ReLU(),
    nn.Linear(256, 10),
)
y = model(x)  # = model[2](model[1](model[0](x)))
\end{lstlisting}

The sub-modules are accessed by integer index or by name (if
constructed from an \code{OrderedDict}).

\subsection{ModuleList}
\label{ssec:module_system:modulelist}

\code{ModuleList} stores a list of sub-modules that the user
applies manually:

\begin{lstlisting}[style=lucid,language=LucidPython]
class StackedLayers(nn.Module):
    def __init__(self, n_layers, dim):
        super().__init__()
        self.layers = nn.ModuleList([
            nn.Linear(dim, dim) for _ in range(n_layers)
        ])

    def forward(self, x):
        for layer in self.layers:
            x = layer(x).relu()
        return x
\end{lstlisting}

Unlike a plain Python list, \code{ModuleList} registers its
contents with the parent module's \code{\_modules} dict, so
\code{model.parameters()} sees every sub-module's parameters.

\subsection{ModuleDict}
\label{ssec:module_system:moduledict}

The dict analogue: \code{ModuleDict} stores named sub-modules. Used
for conditional dispatch, multi-head outputs, mixture-of-experts.

\subsection{ParameterList and ParameterDict}
\label{ssec:module_system:paramlist}

The parameter-level analogues. Useful when a module owns a list of
parameters that are not themselves wrapped in a sub-module.

\section{State Dict and Serialisation}
\label{sec:module_system:state_dict}

\code{module.state\_dict()} returns an ordered dict of every
parameter and buffer in the module tree, with dotted names:

\begin{lstlisting}[style=lucid,language=LucidPython]
sd = model.state_dict()
# {'fc1.weight': tensor(...),
#  'fc1.bias':   tensor(...),
#  'fc2.weight': tensor(...),
#  'fc2.bias':   tensor(...)}

lucid.save(sd, 'checkpoint.pt')

# Later:
sd = lucid.load('checkpoint.pt')
model.load_state_dict(sd)
\end{lstlisting}

\subsection{The strict-loading contract}
\label{ssec:module_system:strict_loading}

\code{load\_state\_dict(sd, strict=True)} checks that every key in
\code{sd} corresponds to a parameter/buffer in the module and that
every parameter/buffer has a key in \code{sd}. A mismatch raises
\code{LucidError}.

\code{strict=False} relaxes the check: missing keys leave their
parameters unchanged, extra keys are silently ignored. Useful for
loading a checkpoint into a partially refactored model.

\subsection{The Bridge boundary}
\label{ssec:module_system:bridge}

\code{state\_dict} touches the numpy bridge boundary
(\chref{ch:bridge_boundaries}) because the persistence format
serialises tensors as numpy arrays. The implementation is in
\code{lucid/serialization/} and \code{lucid/nn/\_state\_dict.py}.

\section{Training and Eval Modes}
\label{sec:module_system:train_eval}

The \code{training} flag toggles behaviour of layers whose forward
differs between training and inference:

\begin{itemize}
  \item \code{Dropout}: zeros out a fraction of activations in
    training, identity in eval.
  \item \code{BatchNorm}: uses batch statistics in training, running
    statistics in eval.
  \item \code{LayerNorm}: unchanged (uses per-sample statistics
    always).
\end{itemize}

The flag is propagated through the module tree by
\code{model.train()} (sets the flag to True) and \code{model.eval()}
(sets to False), recursively.

\subsection{The typical training loop}
\label{ssec:module_system:training_loop}

\begin{lstlisting}[style=lucid,language=LucidPython]
for epoch in range(n_epochs):
    model.train()
    for x, y in train_loader:
        optimizer.zero_grad()
        loss = criterion(model(x), y)
        loss.backward()
        optimizer.step()

    model.eval()
    with lucid.no_grad():
        for x, y in val_loader:
            ...
\end{lstlisting}

The training/eval distinction is the user's responsibility; the
module system provides the toggle but does not enforce its use.

\section{Hooks}
\label{sec:module_system:hooks}

\code{Module} exposes three hook types that the user can attach to
intercept the forward/backward flow:

\begin{itemize}
  \item \textbf{Forward pre-hook}: fires before
    \code{module.forward()}, with the inputs. Can modify the
    inputs.
  \item \textbf{Forward hook}: fires after \code{module.forward()},
    with the inputs and the output. Cannot modify the output (use
    pre-hook plus computed output to modify).
  \item \textbf{Backward hook}: fires during backward, with the
    gradient inputs and outputs.
\end{itemize}

\subsection{Use cases}
\label{ssec:module_system:hook_uses}

\begin{itemize}
  \item Logging activations or gradients during training (for
    debugging vanishing/exploding gradients).
  \item Visualising attention maps (intercept the attention's
    output).
  \item Implementing gradient reversal (multiply by $-1$ in the
    backward hook).
  \item Spectral normalisation (intercept the weight in a forward
    pre-hook, normalise, replace).
\end{itemize}

\subsection{Hook registration}
\label{ssec:module_system:hook_registration}

\begin{lstlisting}[style=lucid,language=LucidPython]
def log_hook(module, inputs, output):
    print(f"{module.__class__.__name__}: out shape {output.shape}")

handle = model.fc1.register_forward_hook(log_hook)

# Run model -- hook fires
y = model(x)

# Remove hook when done
handle.remove()
\end{lstlisting}

The \code{handle.remove()} call is essential: hooks that are not
removed accumulate and slow down every forward pass.

\subsection{Hook firing order}
\label{ssec:module_system:hook_order}

For a single module, the order is:

\begin{enumerate}
  \item Forward pre-hooks fire in registration order;
    each may return a modified \code{inputs} tuple, which is
    threaded to the next pre-hook and finally to
    \code{forward}.
  \item \code{module.forward(*inputs)} runs.
  \item Forward hooks fire in registration order;
    each may inspect (but not modify) the output.
  \item During \code{loss.backward()}, the backward hooks fire
    in reverse order of forward, with the gradient flowing
    backward through the autograd graph.
\end{enumerate}

For a tree of modules, the parent's pre-hook fires before any
child's; the child's full forward (with its own hooks)
completes; then the parent's forward hook fires. The order
mirrors the depth-first walk that the
\code{Module.\_\_call\_\_} performs implicitly.

\subsection{Hook performance cost}
\label{ssec:module_system:hook_perf}

Each hook is a Python function call: $\sim$\,1\,\textmu s
of overhead per call. For modules invoked thousands of times
per iteration (e.g., a deep transformer), the cumulative cost
can become visible.

The typical pattern: attach hooks only during debugging or
data-collection phases, then remove. The
\code{lucid.nn.utils.hook\_scope} context manager makes this
explicit:

\begin{lstlisting}[style=lucid,language=LucidPython]
with lucid.nn.utils.hook_scope() as scope:
    scope.register_forward_hook(model.fc1, log_hook)
    out = model(x)
    # Hooks auto-removed when leaving the with-block
\end{lstlisting}

\subsection{The gradient-reversal hook}
\label{ssec:module_system:grad_reversal}

A well-known use of backward hooks: gradient reversal for
adversarial domain adaptation. The forward is the identity;
the backward multiplies the gradient by $-\lambda$:

\begin{lstlisting}[style=lucid,language=LucidPython]
class GradientReversal(nn.Module):
    def __init__(self, lambda_=1.0):
        super().__init__()
        self.lambda_ = lambda_

    def forward(self, x):
        return x.clone()  # identity

    def __init_subclass__(cls):
        cls.register_backward_hook(
            lambda mod, grad_in, grad_out: (-mod.lambda_ * grad_out[0],)
        )
\end{lstlisting}

A cleaner alternative uses a custom Function
(\chref{ch:custom_functions}); the hook approach is shown for
illustration.

\subsection{State-dict hooks}
\label{ssec:module_system:state_dict_hooks}

Two additional hook families control serialisation:

\begin{itemize}
  \item \code{\_register\_state\_dict\_hook(fn)}: fires after
    \code{state\_dict()} is built, with the resulting dict.
    Can rename keys, add metadata, or strip auxiliary
    tensors before save.
  \item \code{\_register\_load\_state\_dict\_pre\_hook(fn)}:
    fires before \code{load\_state\_dict()} consumes the
    incoming dict. Used by the model-zoo's renaming-table
    machinery (\chref{ch:serialization}~Section~\ref{ssec:serialization:renaming_policy}).
\end{itemize}

\section{Parametrize: Constrained Parameters}
\label{sec:module_system:parametrize}

A frequent advanced pattern: the user wants the model to expose
a logical parameter (say, a weight matrix) that is computed from
an unconstrained underlying parameter (the
\emph{reparameterisation}). The classic example is weight
normalisation \cite{miyato_specnorm_2018}:
\[
  W = g \cdot \frac{V}{\|V\|},
\]
where $V$ is the underlying parameter and $g$ is a learnable
scalar. The optimiser updates $V$ and $g$; the layer sees the
derived $W$.

\Lucid\ provides \code{lucid.nn.utils.parametrize} to register
these reparameterisations cleanly:

\begin{lstlisting}[style=lucid,language=LucidPython]
import lucid.nn.utils.parametrize as P

class WeightNorm(nn.Module):
    def forward(self, weight):
        # weight here is W; returns the reparameterised version
        # Standard weight_norm splits into magnitude g and direction V/||V||
        return weight  # (placeholder; see weight_norm below)

linear = nn.Linear(64, 128)
P.register_parametrization(linear, 'weight', WeightNorm())
# After this, linear.weight goes through WeightNorm.forward on every access.
# The optimiser sees the underlying tensor; linear.weight is a property.
\end{lstlisting}

The mechanism is implemented by replacing the original
parameter with a property that, on read, invokes the
reparameterisation; on write, stores into the underlying
tensor. \code{P.remove\_parametrizations} undoes the
registration, restoring the original parameter (with its
current reparameterised value baked in).

\subsection{Weight normalisation}
\label{ssec:module_system:weight_norm}

The standard wrapper:

\begin{lstlisting}[style=lucid,language=LucidPython]
class _WeightNorm(nn.Module):
    def __init__(self, dim=0):
        super().__init__()
        self.dim = dim
        # After registration the parametrization sees the
        # current weight as `weight`; we split it into magnitude
        # and direction.
        self.g = None   # set on first call
        self.v = None

    def forward(self, weight):
        if self.g is None:
            with lucid.no_grad():
                self.v = nn.Parameter(weight.clone())
                norm = self.v.norm(dim=self.dim, keepdim=True)
                self.g = nn.Parameter(norm)
        return self.g * self.v / self.v.norm(dim=self.dim, keepdim=True)

nn.utils.weight_norm(linear, name='weight', dim=0)  # convenience wrapper
\end{lstlisting}

The split into magnitude $g$ and direction $V / \|V\|$
decouples gradient updates of weight scale from weight
direction; in practice this can improve optimisation on
recurrent and adversarial models.

\subsection{Spectral normalisation}
\label{ssec:module_system:spectral_norm}

Spectral normalisation \cite{miyato_specnorm_2018} bounds the
Lipschitz constant of a layer by dividing the weight matrix by
its largest singular value $\sigma(W)$ (estimated by power
iteration):
\[
  W_{\text{SN}} = W / \sigma(W).
\]
The wrapper maintains two auxiliary vectors $u, v$ as buffers
and updates them in the forward pass via one power iteration:
\[
  v \gets \frac{W^\top u}{\|W^\top u\|}, \qquad
  u \gets \frac{W v}{\|W v\|}, \qquad
  \sigma \approx u^\top W v.
\]
The buffers are updated only in training mode; eval reuses the
last estimate. This stabilises GAN training significantly.

\code{nn.utils.spectral\_norm(layer, name='weight', n\_power\_iterations=1)}
is the standard entry point.

\subsection{Why parametrize, not subclass}
\label{ssec:module_system:why_parametrize}

The same effect could be obtained by subclassing
\code{nn.Linear} and overriding \code{forward}. The advantages
of the parametrize approach:

\begin{itemize}
  \item \textbf{Composability}: multiple reparameterisations
    can be chained on the same parameter
    (\code{register\_parametrization} is idempotent across
    different reparameterisation classes).
  \item \textbf{Transparency}: the original layer's
    \code{forward} is unchanged; the
    reparameterisation is a property modifier.
  \item \textbf{Toggleability}: \code{P.cached()} context
    caches the reparameterised tensor for the duration of a
    forward pass, avoiding re-computation.
\end{itemize}

\section{Apply and Functional Iteration}
\label{sec:module_system:apply}

\code{module.apply(fn)} recursively applies \code{fn} to every
sub-module. Used for initialisation, freezing, and other
tree-wide operations:

\begin{lstlisting}[style=lucid,language=LucidPython]
def init_weights(m):
    if isinstance(m, nn.Linear):
        nn.init.kaiming_normal_(m.weight)
        if m.bias is not None:
            nn.init.zeros_(m.bias)

model.apply(init_weights)
\end{lstlisting}

\subsection{Freezing parameters}
\label{ssec:module_system:freezing}

The standard fine-tuning pattern: freeze the backbone, train only
the head. Implemented by setting \code{requires\_grad = False} on
backbone parameters:

\begin{lstlisting}[style=lucid,language=LucidPython]
for p in model.backbone.parameters():
    p.requires_grad = False

# Optimiser then only sees head parameters
optimizer = lucid.optim.Adam(
    filter(lambda p: p.requires_grad, model.parameters()),
    lr=1e-4,
)
\end{lstlisting}

\section{Moving Modules Across Devices}
\label{sec:module_system:device_move}

\code{module.to(device)} moves every parameter and buffer to the
named device. \code{module.cuda()}, \code{module.cpu()} are
device-specific aliases (\Lucid\ exposes \code{module.metal()} for
the Apple GPU and \code{module.cpu()} for CPU; there is no
\code{cuda()} on Apple silicon).

The move is in-place: existing parameter \code{Tensor} objects are
replaced in the underlying dict. After the move, the optimiser
must be re-bound because it holds references to the old tensors.

\subsection{Casting dtypes}
\label{ssec:module_system:dtype_cast}

\code{module.float()}, \code{module.half()}, \code{module.bfloat16()}
cast every parameter and buffer to the named dtype. The cast
applies recursively.

\section{Cloning, Sharing, and Resetting Modules}
\label{sec:module_system:cloning}

Three operations on an instantiated module are worth treating
carefully because they expose subtleties of the parameter graph.

\subsection{Deep copy}
\label{ssec:module_system:deep_copy}

\code{copy.deepcopy(model)} produces a new module tree with
fresh \code{Parameter} objects (allocated tensors, independent
storage). The autograd graph is not copied; the new parameters
start with no \code{grad\_fn} attached.

\begin{lstlisting}[style=lucid,language=LucidPython]
import copy
model2 = copy.deepcopy(model)
assert model2.fc1.weight is not model.fc1.weight   # different objects
assert (model2.fc1.weight == model.fc1.weight).all()  # same values
\end{lstlisting}

Useful for: snapshot-then-modify experiments, ensemble training
(N copies with different inits), exponential-moving-average
parameter tracking (EMA model is a deepcopy of the live model,
updated each step toward the live parameters).

\subsection{Shared parameters}
\label{ssec:module_system:shared_params}

Assigning the same \code{Parameter} to two different positions
in a tree shares the storage:

\begin{lstlisting}[style=lucid,language=LucidPython]
class TiedEmbedding(nn.Module):
    def __init__(self, vocab_size, dim):
        super().__init__()
        self.embed = nn.Embedding(vocab_size, dim)
        self.head = nn.Linear(dim, vocab_size, bias=False)
        # Tie the weights: head shares storage with embed
        self.head.weight = self.embed.weight
\end{lstlisting}

The optimiser sees both references but updates the underlying
storage once per step. \code{model.parameters()} dedups by
default (de-dupping is via Python's \code{id()} on the underlying
tensor); without dedup, the optimiser would apply the gradient
twice, doubling the effective learning rate for the shared
weights.

Weight tying is the standard idiom in language models
\cite{radford_gpt2_2019}: the input-embedding matrix and the
output-projection matrix are the same tensor, halving the
embedding-parameter count and improving generalisation.

\subsection{Resetting}
\label{ssec:module_system:resetting}

Every layer in \code{lucid.nn} defines \code{reset\_parameters()}
that re-initialises in place. A model-wide reset:

\begin{lstlisting}[style=lucid,language=LucidPython]
def reset_all(module):
    if hasattr(module, 'reset_parameters'):
        module.reset_parameters()

model.apply(reset_all)
\end{lstlisting}

Useful for: starting a new run with the same architecture and
seed, ablation studies, neural-architecture search with shared
weights.

\section{Implementation Notes}
\label{sec:module_system:impl_notes}

The module system is pure Python (no C++ dependency). The base
class is in \code{lucid/nn/module.py}; the parameter class in
\code{lucid/nn/parameter.py}; the containers in
\code{lucid/nn/modules/container.py}.

\subsection{Performance}
\label{ssec:module_system:perf}

The per-call overhead of \code{module(\_)} (the
\code{\_\_call\_\_} dispatch plus hook firing) is a few
microseconds. For modules whose forward kernel takes microseconds
(elementwise small-tensor layers), this can matter; for any
non-trivial forward, it is dominated.

\begin{table}[t]
\centering
\small
\begin{tabular}{lrrr}
\toprule
Module call & Forward kernel ($\mu$s) & \code{\_\_call\_\_} overhead ($\mu$s) & \% overhead \\
\midrule
\code{nn.Linear(16, 16)}    & 12  & 6  & 33\% \\
\code{nn.Linear(256, 256)}  & 35  & 6  & 15\% \\
\code{nn.Linear(4096, 4096)} & 420 & 6 & 1.4\% \\
\code{nn.Conv2d 3x3 (32 ch)} & 95  & 7 & 7\% \\
\code{nn.LayerNorm(768)}     & 18  & 6 & 25\% \\
\code{nn.MultiheadAttention}  & 230 & 8 & 3.4\% \\
\bottomrule
\end{tabular}
\caption[Module call overhead.]{Per-call overhead of
\code{Module.\_\_call\_\_} measured on M3~Max, batch 32.
The overhead is constant (6--8 $\mu$s) regardless of
the underlying kernel; it dominates for tiny layers and
disappears for substantial ones.}
\label{tab:module_system:perf}
\end{table}

The compile path (\chref{ch:compile_design}) eliminates the
per-call overhead inside compiled regions by replacing module
dispatch with direct kernel emission. After
\code{lucid.compile(model)}, even the smallest modules pay zero
Python dispatch cost --- the compiled graph executes as one
MPSGraph submission.

\subsection{The \_\_setattr\_\_ chain}
\label{ssec:module_system:setattr_chain}

The attribute-interception logic at the heart of the module
system:

\begin{lstlisting}[style=lucid,language=LucidPython]
def __setattr__(self, name: str, value):
    if isinstance(value, Parameter):
        self._parameters[name] = value
        # Remove from other dicts if previously assigned
        self._modules.pop(name, None)
        self._buffers.pop(name, None)
    elif isinstance(value, Module):
        self._modules[name] = value
        self._parameters.pop(name, None)
        self._buffers.pop(name, None)
    elif isinstance(value, Tensor) and name in self._buffers:
        # Replacing an existing buffer in-place
        self._buffers[name] = value
    else:
        # Regular Python attribute (numbers, strings, dicts of non-tensors)
        object.__setattr__(self, name, value)
\end{lstlisting}

A subtle case: assigning \code{None} to a previously-registered
parameter slot removes it from \code{\_parameters} but does not
clean up the underlying optimiser state (which was set up at
optimiser construction). The user must re-bind the optimiser
after structural changes to the parameter list.

\subsection{Why a Python base class}
\label{ssec:module_system:why_python}

The module system is in Python rather than C++ for the same reason
the composite layer is in Python
(\chref{ch:composite_ops}~Section~\ref{sec:composite_ops:why}):
the user is most likely to extend modules, and Python is the more
accessible extension language. The cost is the per-call overhead,
which is acceptable for the typical workload.

\section{Summary and Forward References}
\label{sec:module_system:summary}

\code{lucid.nn.Module} is the base class for neural-network
components in \Lucid. Subclasses define \code{\_\_init\_\_} (to
construct sub-modules) and \code{forward} (to apply the
computation); the base class provides parameter collection,
state-dict serialisation, training/eval toggles, device/dtype
movement, and the hook interface. Containers (\code{Sequential},
\code{ModuleList}, \code{ModuleDict}) compose sub-modules.

The next chapter, \chref{ch:linear_conv_norm}, treats the
standard layer library (Linear, Conv, Pool, Norm).
\chref{ch:rnn_family} treats recurrent networks.
\chref{ch:attention} treats attention. \chref{ch:losses} treats
loss functions. \chref{ch:initialization} treats weight
initialisation.

\begin{lucidnote}
For the user: the module system is the API every \refframework\
user already knows. \code{nn.Module}, \code{Parameter},
\code{Sequential}, \code{state\_dict}, \code{train()} /
\code{eval()} --- all work identically. The migration from
\refframework\ to \Lucid\ at the module level is mostly
search-and-replace on the import line.
\end{lucidnote}
